% Shared preamble for all spec documents.
% Include from each spec: % Shared preamble for all spec documents.
% Include from each spec: % Shared preamble for all spec documents.
% Include from each spec: \input{../preamble.tex} (adjust depth).

\documentclass[11pt,a4paper]{article}

\usepackage[utf8]{inputenc}
\usepackage[T1]{fontenc}
\usepackage{lmodern}
\usepackage[a4paper,margin=2.3cm]{geometry}
\usepackage{parskip}
\usepackage{microtype}
\usepackage[dvipsnames]{xcolor}
\usepackage{enumitem}
\usepackage{listings}
\usepackage{tcolorbox}
\usepackage{titlesec}
\usepackage{fancyhdr}
\usepackage{array}
\usepackage{longtable}
\usepackage{booktabs}
\usepackage{amsmath}
\usepackage{amssymb}
\usepackage{hyperref}
\usepackage{cleveref}

\tcbuselibrary{skins,breakable}

\definecolor{specBlue}{RGB}{30,64,132}
\definecolor{specGray}{RGB}{90,90,90}
\definecolor{specAccent}{RGB}{198,40,40}
\definecolor{specGreen}{RGB}{30,110,60}

\hypersetup{
  colorlinks=true,
  linkcolor=specBlue,
  urlcolor=specBlue,
  citecolor=specBlue,
  pdfborder={0 0 0}
}

\titleformat{\section}{\Large\bfseries\color{specBlue}}{\thesection}{0.75em}{}
\titleformat{\subsection}{\large\bfseries\color{specBlue!80}}{\thesubsection}{0.6em}{}

\makeatletter
\newcommand{\specID}[1]{\def\@specID{#1}}
\newcommand{\specTitle}[1]{\def\@specTitle{#1}}
\newcommand{\specStatus}[1]{\def\@specStatus{#1}}
\newcommand{\specVersion}[1]{\def\@specVersion{#1}}
\newcommand{\specOwner}[1]{\def\@specOwner{#1}}
\newcommand{\specTraces}[1]{\def\@specTraces{#1}}

\specID{SPEC-XXX}
\specTitle{Untitled Spec}
\specStatus{DRAFT}
\specVersion{0.1}
\specOwner{Jeremie Nunez}
\specTraces{}

\newcommand{\makespecheader}{%
  \begin{center}
    {\LARGE\bfseries\color{specBlue}\@specTitle}\\[0.4em]
    {\small\color{specGray}\texttt{\@specID} \quad v\@specVersion \quad \textsc{\@specStatus} \quad \@specOwner}
  \end{center}
  \vspace{0.4em}
  \hrule
  \vspace{0.8em}
}

\pagestyle{fancy}
\fancyhf{}
\fancyhead[L]{\small\color{specGray}\@specID}
\fancyhead[R]{\small\color{specGray}\@specTitle}
\fancyfoot[C]{\small\color{specGray}\thepage}
\renewcommand{\headrulewidth}{0pt}
\makeatother

\newtcolorbox{invariant}{
  colback=specBlue!5, colframe=specBlue!75,
  title=Invariant, fonttitle=\bfseries, breakable,
  boxrule=0.5pt, arc=2pt
}

\newtcolorbox{scenario}[1]{
  colback=white, colframe=specBlue!50,
  title={Scenario: #1}, fonttitle=\bfseries, breakable,
  boxrule=0.5pt, arc=2pt
}

\newtcolorbox{ac}[1]{
  colback=specAccent!5, colframe=specAccent!60,
  title={Acceptance Criterion #1}, fonttitle=\bfseries, breakable,
  boxrule=0.5pt, arc=2pt
}

\newtcolorbox{nongoal}{
  colback=specGray!8, colframe=specGray!60,
  title=Non-Goal, fonttitle=\bfseries, breakable,
  boxrule=0.5pt, arc=2pt
}

\lstdefinestyle{codestyle}{
  basicstyle=\ttfamily\small,
  breaklines=true,
  showstringspaces=false,
  frame=single,
  framesep=4pt,
  rulecolor=\color{specBlue!30},
  backgroundcolor=\color{specBlue!2},
  xleftmargin=0.5em,
  xrightmargin=0.5em,
  keywordstyle=\color{specBlue}\bfseries,
  commentstyle=\color{specGray}\itshape,
  stringstyle=\color{specGreen}
}
\lstset{
  inputencoding=utf8,
  extendedchars=true,
  literate=
    {—}{{--}}1
    {–}{{-}}1
    {…}{{...}}1
    {’}{{'}}1
    {‘}{{'}}1
    {“}{{``}}1
    {”}{{''}}1
    {é}{{\'e}}1
    {è}{{\`e}}1
    {ê}{{\^e}}1
    {à}{{\`a}}1
    {â}{{\^a}}1
    {ç}{{\c{c}}}1
    {ô}{{\^o}}1
    {→}{{$\to$}}1
    {←}{{$\leftarrow$}}1
    {×}{{$\times$}}1
    {≥}{{$\geq$}}1
    {≤}{{$\leq$}}1
    {≠}{{$\neq$}}1
    {∈}{{$\in$}}1
    {⌈}{{$\lceil$}}1
    {⌉}{{$\rceil$}}1
    {∞}{{$\infty$}}1
    {α}{{$\alpha$}}1
    {β}{{$\beta$}}1
}
\lstdefinelanguage{TypeScript}{
  keywords={const,let,var,function,return,if,else,for,while,class,interface,type,extends,implements,import,export,from,as,async,await,new,this,true,false,null,undefined,void,any,unknown,never,string,number,boolean,readonly,private,public,protected,enum},
  sensitive=true,
  comment=[l]{//},
  morecomment=[s]{/*}{*/},
  morestring=[b]",
  morestring=[b]',
  morestring=[b]`
}
\lstset{style=codestyle}

\newcommand{\Given}{\textbf{\color{specBlue}Given} }
\newcommand{\When}{\textbf{\color{specBlue}When} }
\newcommand{\Then}{\textbf{\color{specBlue}Then} }
\providecommand{\And}{}
\renewcommand{\And}{\textbf{\color{specBlue}And} }
 (adjust depth).

\documentclass[11pt,a4paper]{article}

\usepackage[utf8]{inputenc}
\usepackage[T1]{fontenc}
\usepackage{lmodern}
\usepackage[a4paper,margin=2.3cm]{geometry}
\usepackage{parskip}
\usepackage{microtype}
\usepackage[dvipsnames]{xcolor}
\usepackage{enumitem}
\usepackage{listings}
\usepackage{tcolorbox}
\usepackage{titlesec}
\usepackage{fancyhdr}
\usepackage{array}
\usepackage{longtable}
\usepackage{booktabs}
\usepackage{amsmath}
\usepackage{amssymb}
\usepackage{hyperref}
\usepackage{cleveref}

\tcbuselibrary{skins,breakable}

\definecolor{specBlue}{RGB}{30,64,132}
\definecolor{specGray}{RGB}{90,90,90}
\definecolor{specAccent}{RGB}{198,40,40}
\definecolor{specGreen}{RGB}{30,110,60}

\hypersetup{
  colorlinks=true,
  linkcolor=specBlue,
  urlcolor=specBlue,
  citecolor=specBlue,
  pdfborder={0 0 0}
}

\titleformat{\section}{\Large\bfseries\color{specBlue}}{\thesection}{0.75em}{}
\titleformat{\subsection}{\large\bfseries\color{specBlue!80}}{\thesubsection}{0.6em}{}

\makeatletter
\newcommand{\specID}[1]{\def\@specID{#1}}
\newcommand{\specTitle}[1]{\def\@specTitle{#1}}
\newcommand{\specStatus}[1]{\def\@specStatus{#1}}
\newcommand{\specVersion}[1]{\def\@specVersion{#1}}
\newcommand{\specOwner}[1]{\def\@specOwner{#1}}
\newcommand{\specTraces}[1]{\def\@specTraces{#1}}

\specID{SPEC-XXX}
\specTitle{Untitled Spec}
\specStatus{DRAFT}
\specVersion{0.1}
\specOwner{Jeremie Nunez}
\specTraces{}

\newcommand{\makespecheader}{%
  \begin{center}
    {\LARGE\bfseries\color{specBlue}\@specTitle}\\[0.4em]
    {\small\color{specGray}\texttt{\@specID} \quad v\@specVersion \quad \textsc{\@specStatus} \quad \@specOwner}
  \end{center}
  \vspace{0.4em}
  \hrule
  \vspace{0.8em}
}

\pagestyle{fancy}
\fancyhf{}
\fancyhead[L]{\small\color{specGray}\@specID}
\fancyhead[R]{\small\color{specGray}\@specTitle}
\fancyfoot[C]{\small\color{specGray}\thepage}
\renewcommand{\headrulewidth}{0pt}
\makeatother

\newtcolorbox{invariant}{
  colback=specBlue!5, colframe=specBlue!75,
  title=Invariant, fonttitle=\bfseries, breakable,
  boxrule=0.5pt, arc=2pt
}

\newtcolorbox{scenario}[1]{
  colback=white, colframe=specBlue!50,
  title={Scenario: #1}, fonttitle=\bfseries, breakable,
  boxrule=0.5pt, arc=2pt
}

\newtcolorbox{ac}[1]{
  colback=specAccent!5, colframe=specAccent!60,
  title={Acceptance Criterion #1}, fonttitle=\bfseries, breakable,
  boxrule=0.5pt, arc=2pt
}

\newtcolorbox{nongoal}{
  colback=specGray!8, colframe=specGray!60,
  title=Non-Goal, fonttitle=\bfseries, breakable,
  boxrule=0.5pt, arc=2pt
}

\lstdefinestyle{codestyle}{
  basicstyle=\ttfamily\small,
  breaklines=true,
  showstringspaces=false,
  frame=single,
  framesep=4pt,
  rulecolor=\color{specBlue!30},
  backgroundcolor=\color{specBlue!2},
  xleftmargin=0.5em,
  xrightmargin=0.5em,
  keywordstyle=\color{specBlue}\bfseries,
  commentstyle=\color{specGray}\itshape,
  stringstyle=\color{specGreen}
}
\lstset{
  inputencoding=utf8,
  extendedchars=true,
  literate=
    {—}{{--}}1
    {–}{{-}}1
    {…}{{...}}1
    {’}{{'}}1
    {‘}{{'}}1
    {“}{{``}}1
    {”}{{''}}1
    {é}{{\'e}}1
    {è}{{\`e}}1
    {ê}{{\^e}}1
    {à}{{\`a}}1
    {â}{{\^a}}1
    {ç}{{\c{c}}}1
    {ô}{{\^o}}1
    {→}{{$\to$}}1
    {←}{{$\leftarrow$}}1
    {×}{{$\times$}}1
    {≥}{{$\geq$}}1
    {≤}{{$\leq$}}1
    {≠}{{$\neq$}}1
    {∈}{{$\in$}}1
    {⌈}{{$\lceil$}}1
    {⌉}{{$\rceil$}}1
    {∞}{{$\infty$}}1
    {α}{{$\alpha$}}1
    {β}{{$\beta$}}1
}
\lstdefinelanguage{TypeScript}{
  keywords={const,let,var,function,return,if,else,for,while,class,interface,type,extends,implements,import,export,from,as,async,await,new,this,true,false,null,undefined,void,any,unknown,never,string,number,boolean,readonly,private,public,protected,enum},
  sensitive=true,
  comment=[l]{//},
  morecomment=[s]{/*}{*/},
  morestring=[b]",
  morestring=[b]',
  morestring=[b]`
}
\lstset{style=codestyle}

\newcommand{\Given}{\textbf{\color{specBlue}Given} }
\newcommand{\When}{\textbf{\color{specBlue}When} }
\newcommand{\Then}{\textbf{\color{specBlue}Then} }
\providecommand{\And}{}
\renewcommand{\And}{\textbf{\color{specBlue}And} }
 (adjust depth).

\documentclass[11pt,a4paper]{article}

\usepackage[utf8]{inputenc}
\usepackage[T1]{fontenc}
\usepackage{lmodern}
\usepackage[a4paper,margin=2.3cm]{geometry}
\usepackage{parskip}
\usepackage{microtype}
\usepackage[dvipsnames]{xcolor}
\usepackage{enumitem}
\usepackage{listings}
\usepackage{tcolorbox}
\usepackage{titlesec}
\usepackage{fancyhdr}
\usepackage{array}
\usepackage{longtable}
\usepackage{booktabs}
\usepackage{amsmath}
\usepackage{amssymb}
\usepackage{hyperref}
\usepackage{cleveref}

\tcbuselibrary{skins,breakable}

\definecolor{specBlue}{RGB}{30,64,132}
\definecolor{specGray}{RGB}{90,90,90}
\definecolor{specAccent}{RGB}{198,40,40}
\definecolor{specGreen}{RGB}{30,110,60}

\hypersetup{
  colorlinks=true,
  linkcolor=specBlue,
  urlcolor=specBlue,
  citecolor=specBlue,
  pdfborder={0 0 0}
}

\titleformat{\section}{\Large\bfseries\color{specBlue}}{\thesection}{0.75em}{}
\titleformat{\subsection}{\large\bfseries\color{specBlue!80}}{\thesubsection}{0.6em}{}

\makeatletter
\newcommand{\specID}[1]{\def\@specID{#1}}
\newcommand{\specTitle}[1]{\def\@specTitle{#1}}
\newcommand{\specStatus}[1]{\def\@specStatus{#1}}
\newcommand{\specVersion}[1]{\def\@specVersion{#1}}
\newcommand{\specOwner}[1]{\def\@specOwner{#1}}
\newcommand{\specTraces}[1]{\def\@specTraces{#1}}

\specID{SPEC-XXX}
\specTitle{Untitled Spec}
\specStatus{DRAFT}
\specVersion{0.1}
\specOwner{Jeremie Nunez}
\specTraces{}

\newcommand{\makespecheader}{%
  \begin{center}
    {\LARGE\bfseries\color{specBlue}\@specTitle}\\[0.4em]
    {\small\color{specGray}\texttt{\@specID} \quad v\@specVersion \quad \textsc{\@specStatus} \quad \@specOwner}
  \end{center}
  \vspace{0.4em}
  \hrule
  \vspace{0.8em}
}

\pagestyle{fancy}
\fancyhf{}
\fancyhead[L]{\small\color{specGray}\@specID}
\fancyhead[R]{\small\color{specGray}\@specTitle}
\fancyfoot[C]{\small\color{specGray}\thepage}
\renewcommand{\headrulewidth}{0pt}
\makeatother

\newtcolorbox{invariant}{
  colback=specBlue!5, colframe=specBlue!75,
  title=Invariant, fonttitle=\bfseries, breakable,
  boxrule=0.5pt, arc=2pt
}

\newtcolorbox{scenario}[1]{
  colback=white, colframe=specBlue!50,
  title={Scenario: #1}, fonttitle=\bfseries, breakable,
  boxrule=0.5pt, arc=2pt
}

\newtcolorbox{ac}[1]{
  colback=specAccent!5, colframe=specAccent!60,
  title={Acceptance Criterion #1}, fonttitle=\bfseries, breakable,
  boxrule=0.5pt, arc=2pt
}

\newtcolorbox{nongoal}{
  colback=specGray!8, colframe=specGray!60,
  title=Non-Goal, fonttitle=\bfseries, breakable,
  boxrule=0.5pt, arc=2pt
}

\lstdefinestyle{codestyle}{
  basicstyle=\ttfamily\small,
  breaklines=true,
  showstringspaces=false,
  frame=single,
  framesep=4pt,
  rulecolor=\color{specBlue!30},
  backgroundcolor=\color{specBlue!2},
  xleftmargin=0.5em,
  xrightmargin=0.5em,
  keywordstyle=\color{specBlue}\bfseries,
  commentstyle=\color{specGray}\itshape,
  stringstyle=\color{specGreen}
}
\lstset{
  inputencoding=utf8,
  extendedchars=true,
  literate=
    {—}{{--}}1
    {–}{{-}}1
    {…}{{...}}1
    {’}{{'}}1
    {‘}{{'}}1
    {“}{{``}}1
    {”}{{''}}1
    {é}{{\'e}}1
    {è}{{\`e}}1
    {ê}{{\^e}}1
    {à}{{\`a}}1
    {â}{{\^a}}1
    {ç}{{\c{c}}}1
    {ô}{{\^o}}1
    {→}{{$\to$}}1
    {←}{{$\leftarrow$}}1
    {×}{{$\times$}}1
    {≥}{{$\geq$}}1
    {≤}{{$\leq$}}1
    {≠}{{$\neq$}}1
    {∈}{{$\in$}}1
    {⌈}{{$\lceil$}}1
    {⌉}{{$\rceil$}}1
    {∞}{{$\infty$}}1
    {α}{{$\alpha$}}1
    {β}{{$\beta$}}1
}
\lstdefinelanguage{TypeScript}{
  keywords={const,let,var,function,return,if,else,for,while,class,interface,type,extends,implements,import,export,from,as,async,await,new,this,true,false,null,undefined,void,any,unknown,never,string,number,boolean,readonly,private,public,protected,enum},
  sensitive=true,
  comment=[l]{//},
  morecomment=[s]{/*}{*/},
  morestring=[b]",
  morestring=[b]',
  morestring=[b]`
}
\lstset{style=codestyle}

\newcommand{\Given}{\textbf{\color{specBlue}Given} }
\newcommand{\When}{\textbf{\color{specBlue}When} }
\newcommand{\Then}{\textbf{\color{specBlue}Then} }
\providecommand{\And}{}
\renewcommand{\And}{\textbf{\color{specBlue}And} }


\specID{SPEC-SW-011}
\specTitle{Editorial Copilot --- structured briefing drafts from audited data}
\specStatus{DRAFT}
\specVersion{0.1}
\specOwner{Jeremie Nunez}

\begin{document}
\makespecheader

\section{Context}
The sweep subsystem produces a stream of audited findings: each finding has been emitted by a
nano worker, classified by category, and (if a reviewer has processed it) either accepted,
rejected, or resolved. The Editorial Copilot turns that curated finding stream into structured
briefing drafts --- short, schema-validated summaries consumable by downstream editors or
dashboards. The point of this service is \emph{not} to be creative: it is to compose audited
inputs into a constrained shape, so that every briefing a human ever reads is traceable back to
an accepted finding, and never to raw model output.

\section{Goals}
\begin{itemize}[leftmargin=1.2em]
  \item Accept a set of finding references and return one briefing object conforming to a strict Zod schema.
  \item Reject any request that would ground the briefing on unaudited content --- only findings with a suitable audit state are eligible.
  \item Validate every model response against the output schema before returning; on failure, retry bounded times, then fail closed.
  \item Expose a deterministic \emph{provenance} mapping: every field in the output cites at least one source finding id.
  \item Run with a bounded token budget: inputs are truncated, outputs are length-capped per field.
\end{itemize}

\section{Non-Goals}
\begin{nongoal}
This spec does not cover: (a) free-form chat or conversational drafting (see SPEC-SW-012),
(b) publication, scheduling, or downstream distribution of briefings, (c) multilingual
translation, (d) the SEO scoring heuristic (owned by the client), (e) human-in-the-loop editing
after a briefing is produced, (f) persistence of briefings beyond returning them to the caller.
\end{nongoal}

\section{Invariants}
\begin{invariant}
No unaudited content is ever emitted. A briefing request whose finding set contains at least one
finding in state \texttt{unaudited} or \texttt{rejected} is refused with a typed error; no model
call is made.
\end{invariant}

\begin{invariant}
The briefing output is structured, not free-form: every response is parsed through the shared
Zod schema \texttt{BriefingSchema} and rejected on failure. The service returns a typed
\texttt{Briefing}, never a raw string.
\end{invariant}

\begin{invariant}
Provenance is total: for every top-level field of a returned briefing, the \texttt{provenance}
map contains at least one source finding id. A briefing with orphan fields is invalid.
\end{invariant}

\begin{invariant}
All text fields are length-bounded at schema level (title $\le$ 120 chars, summary $\le$ 400
chars, each bullet $\le$ 200 chars, $\le$ 6 bullets). Overruns are rejected, not truncated
silently.
\end{invariant}

\begin{invariant}
Model retries on schema-invalid output are bounded by \texttt{MAX\_REPAIR\_ATTEMPTS} (default
1). After the cap, the service returns a typed \texttt{BriefingError} --- never a partially
valid briefing.
\end{invariant}

\section{Interface}
Public surface; consumers include an admin briefing endpoint and internal aggregators.

\begin{lstlisting}[language=TypeScript]
import { z } from "zod";

export const BriefingSchema = z.object({
  title: z.string().min(8).max(120),
  summary: z.string().min(20).max(400),
  bullets: z.array(z.string().min(4).max(200)).min(1).max(6),
  tags: z.array(z.string().min(2).max(40)).max(8),
  category: z.enum([
    "trend", "anomaly", "quality", "coverage", "other",
  ]),
  // Field name -> non-empty list of source finding ids
  provenance: z.record(z.string(), z.array(z.string()).min(1)),
});

export type Briefing = z.infer<typeof BriefingSchema>;

export type BriefingError =
  | { kind: "unaudited_input"; findingIds: string[] }
  | { kind: "schema_invalid"; attempts: number; issues: string[] }
  | { kind: "model_unavailable"; cause: string };

export interface EditorialCopilot {
  draftBriefing(input: {
    findingIds: string[];
    angle?: string;      // optional user-provided framing
  }): Promise<
    | { ok: true; briefing: Briefing }
    | { ok: false; error: BriefingError }
  >;
}
\end{lstlisting}

\section{Scenarios}

\begin{scenario}{Nominal: audited findings produce a valid briefing}
\Given three finding ids all in state \texttt{accepted} \\
\When \texttt{draftBriefing} is invoked with those ids \\
\Then the response is \texttt{\{ ok: true, briefing \}} \\
\And \texttt{briefing} validates against \texttt{BriefingSchema} \\
\And every key of \texttt{briefing.provenance} maps to a non-empty subset of the input ids
\end{scenario}

\begin{scenario}{Refusal on unaudited input}
\Given one of the input finding ids is in state \texttt{unaudited} \\
\When \texttt{draftBriefing} is invoked \\
\Then the service returns \texttt{\{ ok: false, error: \{ kind: "unaudited\_input", findingIds: [\ldots] \} \}} \\
\And no call to the underlying model is made \\
\And the returned \texttt{findingIds} list contains the offending id
\end{scenario}

\begin{scenario}{Schema-invalid model output is repaired once then fails closed}
\Given the model returns a JSON object missing the \texttt{summary} field on attempt 1 \\
\And the model returns a JSON object missing the \texttt{bullets} field on attempt 2 \\
\When \texttt{draftBriefing} is invoked \\
\Then the service returns \texttt{\{ ok: false, error: \{ kind: "schema\_invalid", attempts: 2, issues: [\ldots] \} \}} \\
\And no partially valid briefing is emitted to the caller
\end{scenario}

\begin{scenario}{Orphan field is rejected at validation}
\Given the model returns a briefing with keys \texttt{title, summary, bullets, tags, category, provenance} \\
\And \texttt{provenance} omits the \texttt{summary} key \\
\When validation runs \\
\Then schema validation fails with an \texttt{orphan\_field} issue \\
\And the error propagates as \texttt{schema\_invalid}
\end{scenario}

\begin{scenario}{Model unavailable}
\Given the underlying cortex executor raises a network error \\
\When \texttt{draftBriefing} is invoked \\
\Then the service returns \texttt{\{ ok: false, error: \{ kind: "model\_unavailable", cause \} \}} \\
\And no exception escapes the service boundary
\end{scenario}

\section{Acceptance Criteria}

\begin{ac}{AC-1}
\texttt{draftBriefing} returns \texttt{ok: true} with a \texttt{Briefing} that parses successfully through \texttt{BriefingSchema} when all input findings are audited-accepted.
\end{ac}

\begin{ac}{AC-2}
\texttt{draftBriefing} returns \texttt{unaudited\_input} and performs no model call when at least one input id is \texttt{unaudited} or \texttt{rejected}.
\end{ac}

\begin{ac}{AC-3}
Every top-level field of a returned briefing appears as a key in \texttt{briefing.provenance} with a non-empty list of source ids, all drawn from the input finding ids.
\end{ac}

\begin{ac}{AC-4}
A model response exceeding any schema length bound (e.g. \texttt{summary} $>$ 400 chars) causes the response to be rejected; the caller sees \texttt{schema\_invalid}, never a truncated briefing.
\end{ac}

\begin{ac}{AC-5}
The number of retries on schema failure is capped at \texttt{MAX\_REPAIR\_ATTEMPTS}; beyond the cap the service returns \texttt{schema\_invalid} with the attempt count.
\end{ac}

\begin{ac}{AC-6}
When the cortex executor throws, \texttt{draftBriefing} returns \texttt{model\_unavailable}; no exception reaches the caller.
\end{ac}

\begin{ac}{AC-7}
The service never returns both \texttt{ok: true} and a briefing with an empty \texttt{bullets} array.
\end{ac}

\section{Traceability}

\begin{center}
\begin{tabular}{@{}lll@{}}
\toprule
AC & Test file & Test name \\
\midrule
AC-1 & \texttt{packages/sweep/tests/unit/editorial-copilot.spec.ts} & \texttt{happy path returns schema-valid briefing} \\
AC-2 & \texttt{packages/sweep/tests/unit/editorial-copilot.spec.ts} & \texttt{refuses unaudited findings without model call} \\
AC-3 & \texttt{packages/sweep/tests/unit/editorial-copilot.spec.ts} & \texttt{provenance covers every output field} \\
AC-4 & \texttt{packages/sweep/tests/unit/briefing-schema.spec.ts} & \texttt{rejects overlength fields} \\
AC-5 & \texttt{packages/sweep/tests/integration/editorial-copilot.spec.ts} & \texttt{schema repair retries are capped} \\
AC-6 & \texttt{packages/sweep/tests/integration/editorial-copilot.spec.ts} & \texttt{model failure returns model\_unavailable} \\
AC-7 & \texttt{packages/sweep/tests/unit/briefing-schema.spec.ts} & \texttt{bullets non-empty invariant} \\
\bottomrule
\end{tabular}
\end{center}

\section{Open Questions}
\begin{itemize}[leftmargin=1.2em]
  \item Should \texttt{category} be inferred from the dominant finding category, or selected freely by the model within the enum?
  \item Do we want a soft retry on \texttt{model\_unavailable} (with backoff), or leave retry policy to the caller?
  \item Should we expose a dry-run mode that returns only the prompt that would be sent, for audit purposes?
  \item How do we version \texttt{BriefingSchema}? A breaking field change affects all existing consumers.
\end{itemize}

\end{document}
